\documentclass[11pt]{article}
\usepackage[a4paper,margin=25mm]{geometry}
\usepackage{amsmath,amssymb,amsthm,booktabs,tikz}
\usepackage[colorlinks=true,urlcolor=blue,citecolor=blue]{hyperref}
\hypersetup{pdftitle={The sharp corner threshold for five-colour cubes},pdfauthor={Jason Davies}}
\newtheorem{theorem}{Theorem}
\newtheorem{lemma}[theorem]{Lemma}
\newtheorem{proposition}[theorem]{Proposition}
\theoremstyle{remark}
\newtheorem*{remark}{Remark}
\newcommand{\col}{\{0,1,2,3,4\}}
\newcommand{\mass}[2]{#1(#2)}
\title{The sharp corner threshold for five-colour cubes}
\author{Jason Davies}
\date{November 3, 2026}
\begin{document}
\maketitle
\begin{abstract}
Every collection of fifteen cubes, each displaying all five colours of a
fixed palette, contains eight that form a $2\times2\times2$ cube with monochromatic
exterior faces using all five colours. Fourteen need not suffice. This
settles the five-colour corner value listed as unknown by Berkove,
Eickemeyer and Schulman (2018). An elementary Hall argument then shows that
any $n^3$ such cubes can all form a complete $n\times n\times n$ cube for
$n\ge3$, and that $12n-16$ cubes suffice for its frame when $n\ge10$.
With at most $k$ source varieties, the sharp corner threshold is $8$ for
$k=1$, $13$ for $k=2$, and $15$ for $k\ge3$. The corner lower bound is
elementary; the upper bound uses 279 finite interval certificates. Lean
verifies these certificates, their connection to physical arrangements,
and the deductions for larger cubes and bounded numbers of varieties.
\end{abstract}
\noindent\small\textit{AI transparency.} The research, code and manuscript were
developed using GPT-6-Astra (OpenAI), through Codex, under the author's
direction. This working manuscript has not received independent human peer
review; Section~\ref{sec:formal} states the formal verification scope.
\normalsize

\section{Introduction}
MacMahon's coloured-cube puzzles~\cite{MacMahon} ask when a collection of
face-coloured cubes can form a larger cube with monochromatic exterior
faces. Here every source and target uses all five colours of a fixed
palette. Proper rotations identify varieties; reflections are not allowed
as moves. Collections may contain repeated varieties, and touching interior
faces need not have equal colours. The target may depend on the collection.

For $n\ge2$, an \emph{$n$-frame} consists of the corners and edges of an
$n\times n\times n$ cube, using $12n-16$ pieces. Write $F_5(n)$ for the
least available pool size guaranteeing an $n$-frame, and $G_5(n)$ for the
least pool size guaranteeing a complete $n$-cube. An assembly may use a
subset of the pool. At $n=2$, the frame and complete cube both consist of
eight corners; we write $N_5=F_5(2)=G_5(2)$.

Berkove, Eickemeyer and Schulman~\cite[Section 8, Table 7]{BES} tabulate
frame numbers for palettes of two to six colours, leaving the five-colour
column unknown. Their corner values for two, three, four and six colours
are respectively $11,23,11,24$. We determine the missing value.
\begin{theorem}\label{thm:main}
The sharp five-colour corner threshold is $N_5=15$.
\end{theorem}
The conclusion for larger collections follows by choosing any fifteen cubes.
A short corner argument gives an explicit fourteen-cube obstruction.
For the upper bound, Hall deficiencies and symmetry reduce hypothetical
counterexamples to 357 finite cases. Containment, budget bounds and a
deletion lemma leave 279 smaller problems, each excluded by an interval
certificate (Sections~\ref{sec:cases}--\ref{sec:certificates}).

The corner theorem also yields
\[
 G_5(n)=\max\{n^3,15\}\qquad(n\ge2),
\]
and the frame bounds
\[
 12n-16\le F_5(n)\le\max\{12n-16,11n-6\}\qquad(n\ge3).
\]
Thus every 27 cubes form a complete $3$-cube without spares, and
$F_5(n)=12n-16$ for $n\ge10$. The argument in Section~\ref{sec:frames}
uses only the corner theorem, elementary geometry and Hall's theorem;
it requires no further interval certificates. Section~\ref{sec:support}
determines the thresholds with a bounded number of source varieties,
showing in particular that a fourteen-cube obstruction needs at least
three varieties. All these results are formalised in Lean.
The accompanying review archive contains the certificates, programs,
geometric audits and Lean sources.\footnote{Intended repository:
\url{https://github.com/jasondavies/five-colour-cubes}.
Public release is pending; the review archive includes source checksums,
the base Git revision and bundled pinned project dependencies.}

The companion paper~\cite{FourToThree} considers four-colour sources and
targets using exactly three colours. It gives the sharp frame threshold
$16n-3$ for every side length $n\ge2$, with a classification of the largest
failing collections. The two developments share the palette-independent
geometry in \texttt{ColourCubes}~\cite{ColourCubes};
neither threshold theorem is used to prove the other.

A limited literature search has found no later value for the five-colour entry.
This is not a comprehensive priority claim. In particular, Rosenfeld's
October 2026 article~\cite{Rosenfeld} has not been examined in full.

\section{Geometry and Hall deficiencies}\label{sec:model}
Represent a cube by its colours on the faces
$(+x,-x,+y,-y,+z,-z)$ in a right-handed coordinate frame. A five-colour cube
has one repeated colour. There are 75 varieties: for each repeated colour,
three have its two faces opposite and twelve have them adjacent. In the
opposite case, the other four colours form two opposite pairs, giving three
possibilities. In the adjacent case, fix the two equal faces; the 24 orders
of the other colours are identified in pairs by the rotation exchanging
those equal faces. Thus there are 15 opposite-repeat varieties and 60
adjacent-repeat varieties. Each class is one orbit under colour permutations.

A corner type is the ordered triple of its incident face colours, modulo
cyclic permutation. Order the outward face normals positively, equivalently
read the triple counterclockwise from outside along the body diagonal.
This reverses the clockwise convention in~\cite{BES}, consistently on every
cube, and therefore preserves all matching and containment relations.
Repeated occurrences of a corner type are retained as distinct target
positions.

\begin{lemma}[Corner--rotation correspondence]\label{lem:corner}
A source cube fits a target corner by a proper rotation if and only if it
has the same cyclic corner type at some corner.
\end{lemma}
\begin{proof}
A proper rotation preserves the orientation of a corner frame, so it induces
a cyclic permutation of its three outward faces. Conversely, any two
positively oriented corner frames, with a prescribed cyclic alignment,
are related by a determinant-positive signed permutation of the coordinate
axes. This is a proper rotation of the cube.
\end{proof}
The palette-independent Lean version of this lemma is reused from the
shared \texttt{ColourCubes} library~\cite{ColourCubes}.

Fix a target $T$. Form a bipartite graph whose left vertices are its eight
corner positions and whose right vertices are the individual source copies.
A copy is adjacent to a position when it fits there. A matching covering
the eight positions is exactly a solution: each chosen copy can be rotated
independently, since interior faces are unrestricted. By Hall's
theorem~\cite{Hall}, failure is equivalent to a nonempty set $D$ of target
positions having fewer than $|D|$ neighbouring copies.

Index the varieties by $i=1,\ldots,75$, and let $x_i$ be their multiplicities.
For a set $S$ of varieties, write $x(S)=\sum_{i\in S}x_i$.
Let $N_T(D)$ be the varieties fitting at least one position of $D$. Then
failure of $T$ is equivalent to the existence of a cut
\begin{equation}\label{eq:hall}
  x(N_T(D))\le |D|-1,\qquad \varnothing\ne D\subseteq\{1,\ldots,8\}.
\end{equation}
There are 255 cuts per target, hence 19,125 cuts altogether.

\section{An elementary fourteen-cube obstruction}\label{sec:lower}
Consider four copies of $A$, five of $B$, and five of $C$, where
\[
 A=(0,0,1,2,3,4),\quad B=(0,1,2,3,2,4),\quad C=(0,3,1,2,4,1).
\]
Their corner types are shown below. Triples denote cyclic classes; the last
column counts positions whose type occurs in neither of the other varieties.
\begin{center}
\begin{tabular}{clc}
\toprule
Variety & Corner types & Private positions\\
\midrule
$A$ & $013,014,023,024,031,032,041,042$ & 5\\
$B$ & $022,023,034,042,122,124,132,143$ & 6\\
$C$ & $011,014,021,042,113,123,134,243$ & 6\\
\bottomrule
\end{tabular}
\end{center}
\begin{lemma}\label{lem:lower}
The collection $4A+5B+5C$ has no five-colour solution.
\end{lemma}
\begin{proof}
A solvable target must have every corner type in the displayed union.
Its only two-colour types are $022,122,011,113$. An adjacent-repeat target
therefore repeats either colour 2, with types $022,122$, or colour 1,
with types $011,113$. In either case these two corners determine the colours
on the faces at the ends of the repeated-colour edge. The two orders of the
remaining colours give a mirror pair: $B$ and its mirror, or $C$ and its
mirror. The mirror of $B$ requires $043$, and the mirror of $C$ requires
$012$; both types are absent. Only $B,C$ remain.

An opposite-repeat target requires both cyclic orientations of each corner
triple. The unordered triples with both orientations in the union are
\[
  013,014,023,024,123,134.
\]
If its repeated colour is $r$, delete $r$ from the four triples containing
it. The remaining pairs must form a $K_{2,2}$ on the other four colours:
choose one face from each of their two opposite pairs. Colours 2 and 4
occur in fewer than four listed triples. For $r=1$ or $r=3$, the four pairs
have degree sequence $(3,2,2,1)$, so do not form $K_{2,2}$. For $r=0$, the
pairs are $13,14,23,24$, forcing opposite pairs $\{1,2\}$ and $\{3,4\}$.
This is $A$.

Thus only $A,B,C$ can be targets. Their private positions require respectively
at least five, six and six copies of the corresponding variety, exceeding
the available four, five and five copies (Figure~\ref{fig:obstruction}).
None is solvable.
\end{proof}

\begin{figure}[ht]
\centering
% Generated by scripts/generate_five_colour_figure.py; do not edit.
\definecolor{cube0}{HTML}{E69F00}
\definecolor{cube1}{HTML}{56B4E9}
\definecolor{cube2}{HTML}{009E73}
\definecolor{cube3}{HTML}{F0E442}
\definecolor{cube4}{HTML}{CC79A7}
\begin{tikzpicture}[x=0.198mm,y=-0.198mm]
\path[use as bounding box] (0,0) rectangle (800,470);
\node[anchor=center,inner sep=0pt,outer sep=0pt,text=black!93,font={\sffamily\fontsize{10}{12}\selectfont\bfseries}] at (72,22) {Collection};
\node[anchor=center,inner sep=0pt,outer sep=0pt,text=black!93,font={\sffamily\fontsize{10}{12}\selectfont\bfseries}] at (372,22) {Corner types of each variety};
\node[anchor=west,inner sep=0pt,outer sep=0pt,text=black!93,font={\sffamily\fontsize{10}{12}\selectfont\bfseries}] at (604,22) {Why it fails};
\draw[draw=black!20,line width=0.674pt] (16,40) -- (784,40);
\draw[fill=cube3,draw=black!87,line join=round,line width=0.674pt] (48,54) -- (72,54) -- (72,78) -- (48,78) -- cycle;
\node[anchor=center,inner sep=0pt,outer sep=0pt,text=black!93,font={\sffamily\fontsize{10}{12}\selectfont}] at (60,66) {3};
\draw[fill=cube2,draw=black!87,line join=round,line width=0.674pt] (24,78) -- (48,78) -- (48,102) -- (24,102) -- cycle;
\node[anchor=center,inner sep=0pt,outer sep=0pt,text=black!93,font={\sffamily\fontsize{10}{12}\selectfont}] at (36,90) {2};
\draw[fill=cube0,draw=black!87,line join=round,line width=0.674pt] (48,78) -- (72,78) -- (72,102) -- (48,102) -- cycle;
\node[anchor=center,inner sep=0pt,outer sep=0pt,text=black!93,font={\sffamily\fontsize{10}{12}\selectfont}] at (60,90) {0};
\draw[fill=cube1,draw=black!87,line join=round,line width=0.674pt] (72,78) -- (96,78) -- (96,102) -- (72,102) -- cycle;
\node[anchor=center,inner sep=0pt,outer sep=0pt,text=black!93,font={\sffamily\fontsize{10}{12}\selectfont}] at (84,90) {1};
\draw[fill=cube0,draw=black!87,line join=round,line width=0.674pt] (96,78) -- (120,78) -- (120,102) -- (96,102) -- cycle;
\node[anchor=center,inner sep=0pt,outer sep=0pt,text=black!93,font={\sffamily\fontsize{10}{12}\selectfont}] at (108,90) {0};
\draw[fill=cube4,draw=black!87,line join=round,line width=0.674pt] (48,102) -- (72,102) -- (72,126) -- (48,126) -- cycle;
\node[anchor=center,inner sep=0pt,outer sep=0pt,text=black!93,font={\sffamily\fontsize{10}{12}\selectfont}] at (60,114) {4};
\node[anchor=center,inner sep=0pt,outer sep=0pt,text=black!93,font={\sffamily\fontsize{10}{12}\selectfont\bfseries}] at (72,148) {4 \(\times\) A};
\node[anchor=center,inner sep=0pt,outer sep=0pt,text=black!93,font={\sffamily\fontsize{9}{10.8}\selectfont}] at (72,170) {001234};
\draw[fill=cube0,draw=black!87,line join=round,line width=0.674pt] (180,96) -- (162.679,86) -- (162.679,106) -- (180,116) -- cycle;
\node[anchor=center,inner sep=0pt,outer sep=0pt,text=black!93,font={\sffamily\fontsize{8}{9.6}\selectfont}] at (171.34,101) {0};
\draw[fill=cube1,draw=black!87,line join=round,line width=0.674pt] (180,96) -- (180,116) -- (197.321,106) -- (197.321,86) -- cycle;
\node[anchor=center,inner sep=0pt,outer sep=0pt,text=black!93,font={\sffamily\fontsize{8}{9.6}\selectfont}] at (188.66,101) {1};
\draw[fill=cube3,draw=black!87,line join=round,line width=0.674pt] (180,96) -- (197.321,86) -- (180,76) -- (162.679,86) -- cycle;
\node[anchor=center,inner sep=0pt,outer sep=0pt,text=black!93,font={\sffamily\fontsize{8}{9.6}\selectfont}] at (180,86) {3};
\node[anchor=center,inner sep=0pt,outer sep=0pt,text=black!93,font={\sffamily\fontsize{9}{10.8}\selectfont}] at (180,131) {013};
\draw[fill=cube0,draw=black!87,line join=round,line width=0.674pt] (234,96) -- (216.679,86) -- (216.679,106) -- (234,116) -- cycle;
\node[anchor=center,inner sep=0pt,outer sep=0pt,text=black!93,font={\sffamily\fontsize{8}{9.6}\selectfont}] at (225.34,101) {0};
\draw[fill=cube2,draw=black!87,line join=round,line width=0.674pt] (234,96) -- (234,116) -- (251.321,106) -- (251.321,86) -- cycle;
\node[anchor=center,inner sep=0pt,outer sep=0pt,text=black!93,font={\sffamily\fontsize{8}{9.6}\selectfont}] at (242.66,101) {2};
\draw[fill=cube4,draw=black!87,line join=round,line width=0.674pt] (234,96) -- (251.321,86) -- (234,76) -- (216.679,86) -- cycle;
\node[anchor=center,inner sep=0pt,outer sep=0pt,text=black!93,font={\sffamily\fontsize{8}{9.6}\selectfont}] at (234,86) {4};
\node[anchor=center,inner sep=0pt,outer sep=0pt,text=black!93,font={\sffamily\fontsize{9}{10.8}\selectfont}] at (234,131) {024};
\draw[fill=cube0,draw=black!87,line join=round,line width=0.674pt] (288,96) -- (270.679,86) -- (270.679,106) -- (288,116) -- cycle;
\node[anchor=center,inner sep=0pt,outer sep=0pt,text=black!93,font={\sffamily\fontsize{8}{9.6}\selectfont}] at (279.34,101) {0};
\draw[fill=cube3,draw=black!87,line join=round,line width=0.674pt] (288,96) -- (288,116) -- (305.321,106) -- (305.321,86) -- cycle;
\node[anchor=center,inner sep=0pt,outer sep=0pt,text=black!93,font={\sffamily\fontsize{8}{9.6}\selectfont}] at (296.66,101) {3};
\draw[fill=cube1,draw=black!87,line join=round,line width=0.674pt] (288,96) -- (305.321,86) -- (288,76) -- (270.679,86) -- cycle;
\node[anchor=center,inner sep=0pt,outer sep=0pt,text=black!93,font={\sffamily\fontsize{8}{9.6}\selectfont}] at (288,86) {1};
\node[anchor=center,inner sep=0pt,outer sep=0pt,text=black!93,font={\sffamily\fontsize{9}{10.8}\selectfont}] at (288,131) {031};
\draw[fill=cube0,draw=black!87,line join=round,line width=0.674pt] (342,96) -- (324.679,86) -- (324.679,106) -- (342,116) -- cycle;
\node[anchor=center,inner sep=0pt,outer sep=0pt,text=black!93,font={\sffamily\fontsize{8}{9.6}\selectfont}] at (333.34,101) {0};
\draw[fill=cube3,draw=black!87,line join=round,line width=0.674pt] (342,96) -- (342,116) -- (359.321,106) -- (359.321,86) -- cycle;
\node[anchor=center,inner sep=0pt,outer sep=0pt,text=black!93,font={\sffamily\fontsize{8}{9.6}\selectfont}] at (350.66,101) {3};
\draw[fill=cube2,draw=black!87,line join=round,line width=0.674pt] (342,96) -- (359.321,86) -- (342,76) -- (324.679,86) -- cycle;
\node[anchor=center,inner sep=0pt,outer sep=0pt,text=black!93,font={\sffamily\fontsize{8}{9.6}\selectfont}] at (342,86) {2};
\node[anchor=center,inner sep=0pt,outer sep=0pt,text=black!93,font={\sffamily\fontsize{9}{10.8}\selectfont}] at (342,131) {032};
\draw[fill=cube0,draw=black!87,line join=round,line width=0.674pt] (396,96) -- (378.679,86) -- (378.679,106) -- (396,116) -- cycle;
\node[anchor=center,inner sep=0pt,outer sep=0pt,text=black!93,font={\sffamily\fontsize{8}{9.6}\selectfont}] at (387.34,101) {0};
\draw[fill=cube4,draw=black!87,line join=round,line width=0.674pt] (396,96) -- (396,116) -- (413.321,106) -- (413.321,86) -- cycle;
\node[anchor=center,inner sep=0pt,outer sep=0pt,text=black!93,font={\sffamily\fontsize{8}{9.6}\selectfont}] at (404.66,101) {4};
\draw[fill=cube1,draw=black!87,line join=round,line width=0.674pt] (396,96) -- (413.321,86) -- (396,76) -- (378.679,86) -- cycle;
\node[anchor=center,inner sep=0pt,outer sep=0pt,text=black!93,font={\sffamily\fontsize{8}{9.6}\selectfont}] at (396,86) {1};
\node[anchor=center,inner sep=0pt,outer sep=0pt,text=black!93,font={\sffamily\fontsize{9}{10.8}\selectfont}] at (396,131) {041};
\draw[fill=cube0!30,draw=black!40,line join=round,line width=0.674pt] (450,96) -- (432.679,86) -- (432.679,106) -- (450,116) -- cycle;
\node[anchor=center,inner sep=0pt,outer sep=0pt,text=black!60,font={\sffamily\fontsize{8}{9.6}\selectfont}] at (441.34,101) {0};
\draw[fill=cube1!30,draw=black!40,line join=round,line width=0.674pt] (450,96) -- (450,116) -- (467.321,106) -- (467.321,86) -- cycle;
\node[anchor=center,inner sep=0pt,outer sep=0pt,text=black!60,font={\sffamily\fontsize{8}{9.6}\selectfont}] at (458.66,101) {1};
\draw[fill=cube4!30,draw=black!40,line join=round,line width=0.674pt] (450,96) -- (467.321,86) -- (450,76) -- (432.679,86) -- cycle;
\node[anchor=center,inner sep=0pt,outer sep=0pt,text=black!60,font={\sffamily\fontsize{8}{9.6}\selectfont}] at (450,86) {4};
\node[anchor=center,inner sep=0pt,outer sep=0pt,text=black!60,font={\sffamily\fontsize{9}{10.8}\selectfont}] at (450,131) {014};
\draw[fill=cube0!30,draw=black!40,line join=round,line width=0.674pt] (504,96) -- (486.679,86) -- (486.679,106) -- (504,116) -- cycle;
\node[anchor=center,inner sep=0pt,outer sep=0pt,text=black!60,font={\sffamily\fontsize{8}{9.6}\selectfont}] at (495.34,101) {0};
\draw[fill=cube2!30,draw=black!40,line join=round,line width=0.674pt] (504,96) -- (504,116) -- (521.321,106) -- (521.321,86) -- cycle;
\node[anchor=center,inner sep=0pt,outer sep=0pt,text=black!60,font={\sffamily\fontsize{8}{9.6}\selectfont}] at (512.66,101) {2};
\draw[fill=cube3!30,draw=black!40,line join=round,line width=0.674pt] (504,96) -- (521.321,86) -- (504,76) -- (486.679,86) -- cycle;
\node[anchor=center,inner sep=0pt,outer sep=0pt,text=black!60,font={\sffamily\fontsize{8}{9.6}\selectfont}] at (504,86) {3};
\node[anchor=center,inner sep=0pt,outer sep=0pt,text=black!60,font={\sffamily\fontsize{9}{10.8}\selectfont}] at (504,131) {023};
\draw[fill=cube0!30,draw=black!40,line join=round,line width=0.674pt] (558,96) -- (540.679,86) -- (540.679,106) -- (558,116) -- cycle;
\node[anchor=center,inner sep=0pt,outer sep=0pt,text=black!60,font={\sffamily\fontsize{8}{9.6}\selectfont}] at (549.34,101) {0};
\draw[fill=cube4!30,draw=black!40,line join=round,line width=0.674pt] (558,96) -- (558,116) -- (575.321,106) -- (575.321,86) -- cycle;
\node[anchor=center,inner sep=0pt,outer sep=0pt,text=black!60,font={\sffamily\fontsize{8}{9.6}\selectfont}] at (566.66,101) {4};
\draw[fill=cube2!30,draw=black!40,line join=round,line width=0.674pt] (558,96) -- (575.321,86) -- (558,76) -- (540.679,86) -- cycle;
\node[anchor=center,inner sep=0pt,outer sep=0pt,text=black!60,font={\sffamily\fontsize{8}{9.6}\selectfont}] at (558,86) {2};
\node[anchor=center,inner sep=0pt,outer sep=0pt,text=black!60,font={\sffamily\fontsize{9}{10.8}\selectfont}] at (558,131) {042};
\draw[draw=black!87,line width=0.787pt] (158,144) -- (158,150) -- (418,150) -- (418,144);
\node[anchor=center,inner sep=0pt,outer sep=0pt,text=black!93,font={\sffamily\fontsize{9}{10.8}\selectfont\bfseries}] at (288,168) {5 private};
\node[anchor=west,inner sep=0pt,outer sep=0pt,text=black!93,font={\sffamily\fontsize{10}{12}\selectfont}] at (604,86) {5 corners need A;};
\node[anchor=west,inner sep=0pt,outer sep=0pt,text=black!93,font={\sffamily\fontsize{10}{12}\selectfont\bfseries}] at (604,112) {only 4 copies.};
\draw[draw=black!20,line width=0.674pt] (16,180) -- (784,180);
\draw[fill=cube2,draw=black!87,line join=round,line width=0.674pt] (48,194) -- (72,194) -- (72,218) -- (48,218) -- cycle;
\node[anchor=center,inner sep=0pt,outer sep=0pt,text=black!93,font={\sffamily\fontsize{10}{12}\selectfont}] at (60,206) {2};
\draw[fill=cube3,draw=black!87,line join=round,line width=0.674pt] (24,218) -- (48,218) -- (48,242) -- (24,242) -- cycle;
\node[anchor=center,inner sep=0pt,outer sep=0pt,text=black!93,font={\sffamily\fontsize{10}{12}\selectfont}] at (36,230) {3};
\draw[fill=cube0,draw=black!87,line join=round,line width=0.674pt] (48,218) -- (72,218) -- (72,242) -- (48,242) -- cycle;
\node[anchor=center,inner sep=0pt,outer sep=0pt,text=black!93,font={\sffamily\fontsize{10}{12}\selectfont}] at (60,230) {0};
\draw[fill=cube2,draw=black!87,line join=round,line width=0.674pt] (72,218) -- (96,218) -- (96,242) -- (72,242) -- cycle;
\node[anchor=center,inner sep=0pt,outer sep=0pt,text=black!93,font={\sffamily\fontsize{10}{12}\selectfont}] at (84,230) {2};
\draw[fill=cube1,draw=black!87,line join=round,line width=0.674pt] (96,218) -- (120,218) -- (120,242) -- (96,242) -- cycle;
\node[anchor=center,inner sep=0pt,outer sep=0pt,text=black!93,font={\sffamily\fontsize{10}{12}\selectfont}] at (108,230) {1};
\draw[fill=cube4,draw=black!87,line join=round,line width=0.674pt] (48,242) -- (72,242) -- (72,266) -- (48,266) -- cycle;
\node[anchor=center,inner sep=0pt,outer sep=0pt,text=black!93,font={\sffamily\fontsize{10}{12}\selectfont}] at (60,254) {4};
\node[anchor=center,inner sep=0pt,outer sep=0pt,text=black!93,font={\sffamily\fontsize{10}{12}\selectfont\bfseries}] at (72,288) {5 \(\times\) B};
\node[anchor=center,inner sep=0pt,outer sep=0pt,text=black!93,font={\sffamily\fontsize{9}{10.8}\selectfont}] at (72,310) {012324};
\draw[fill=cube0,draw=black!87,line join=round,line width=0.674pt] (180,236) -- (162.679,226) -- (162.679,246) -- (180,256) -- cycle;
\node[anchor=center,inner sep=0pt,outer sep=0pt,text=black!93,font={\sffamily\fontsize{8}{9.6}\selectfont}] at (171.34,241) {0};
\draw[fill=cube2,draw=black!87,line join=round,line width=0.674pt] (180,236) -- (180,256) -- (197.321,246) -- (197.321,226) -- cycle;
\node[anchor=center,inner sep=0pt,outer sep=0pt,text=black!93,font={\sffamily\fontsize{8}{9.6}\selectfont}] at (188.66,241) {2};
\draw[fill=cube2,draw=black!87,line join=round,line width=0.674pt] (180,236) -- (197.321,226) -- (180,216) -- (162.679,226) -- cycle;
\node[anchor=center,inner sep=0pt,outer sep=0pt,text=black!93,font={\sffamily\fontsize{8}{9.6}\selectfont}] at (180,226) {2};
\node[anchor=center,inner sep=0pt,outer sep=0pt,text=black!93,font={\sffamily\fontsize{9}{10.8}\selectfont}] at (180,271) {022};
\draw[fill=cube0,draw=black!87,line join=round,line width=0.674pt] (234,236) -- (216.679,226) -- (216.679,246) -- (234,256) -- cycle;
\node[anchor=center,inner sep=0pt,outer sep=0pt,text=black!93,font={\sffamily\fontsize{8}{9.6}\selectfont}] at (225.34,241) {0};
\draw[fill=cube3,draw=black!87,line join=round,line width=0.674pt] (234,236) -- (234,256) -- (251.321,246) -- (251.321,226) -- cycle;
\node[anchor=center,inner sep=0pt,outer sep=0pt,text=black!93,font={\sffamily\fontsize{8}{9.6}\selectfont}] at (242.66,241) {3};
\draw[fill=cube4,draw=black!87,line join=round,line width=0.674pt] (234,236) -- (251.321,226) -- (234,216) -- (216.679,226) -- cycle;
\node[anchor=center,inner sep=0pt,outer sep=0pt,text=black!93,font={\sffamily\fontsize{8}{9.6}\selectfont}] at (234,226) {4};
\node[anchor=center,inner sep=0pt,outer sep=0pt,text=black!93,font={\sffamily\fontsize{9}{10.8}\selectfont}] at (234,271) {034};
\draw[fill=cube1,draw=black!87,line join=round,line width=0.674pt] (288,236) -- (270.679,226) -- (270.679,246) -- (288,256) -- cycle;
\node[anchor=center,inner sep=0pt,outer sep=0pt,text=black!93,font={\sffamily\fontsize{8}{9.6}\selectfont}] at (279.34,241) {1};
\draw[fill=cube2,draw=black!87,line join=round,line width=0.674pt] (288,236) -- (288,256) -- (305.321,246) -- (305.321,226) -- cycle;
\node[anchor=center,inner sep=0pt,outer sep=0pt,text=black!93,font={\sffamily\fontsize{8}{9.6}\selectfont}] at (296.66,241) {2};
\draw[fill=cube2,draw=black!87,line join=round,line width=0.674pt] (288,236) -- (305.321,226) -- (288,216) -- (270.679,226) -- cycle;
\node[anchor=center,inner sep=0pt,outer sep=0pt,text=black!93,font={\sffamily\fontsize{8}{9.6}\selectfont}] at (288,226) {2};
\node[anchor=center,inner sep=0pt,outer sep=0pt,text=black!93,font={\sffamily\fontsize{9}{10.8}\selectfont}] at (288,271) {122};
\draw[fill=cube1,draw=black!87,line join=round,line width=0.674pt] (342,236) -- (324.679,226) -- (324.679,246) -- (342,256) -- cycle;
\node[anchor=center,inner sep=0pt,outer sep=0pt,text=black!93,font={\sffamily\fontsize{8}{9.6}\selectfont}] at (333.34,241) {1};
\draw[fill=cube2,draw=black!87,line join=round,line width=0.674pt] (342,236) -- (342,256) -- (359.321,246) -- (359.321,226) -- cycle;
\node[anchor=center,inner sep=0pt,outer sep=0pt,text=black!93,font={\sffamily\fontsize{8}{9.6}\selectfont}] at (350.66,241) {2};
\draw[fill=cube4,draw=black!87,line join=round,line width=0.674pt] (342,236) -- (359.321,226) -- (342,216) -- (324.679,226) -- cycle;
\node[anchor=center,inner sep=0pt,outer sep=0pt,text=black!93,font={\sffamily\fontsize{8}{9.6}\selectfont}] at (342,226) {4};
\node[anchor=center,inner sep=0pt,outer sep=0pt,text=black!93,font={\sffamily\fontsize{9}{10.8}\selectfont}] at (342,271) {124};
\draw[fill=cube1,draw=black!87,line join=round,line width=0.674pt] (396,236) -- (378.679,226) -- (378.679,246) -- (396,256) -- cycle;
\node[anchor=center,inner sep=0pt,outer sep=0pt,text=black!93,font={\sffamily\fontsize{8}{9.6}\selectfont}] at (387.34,241) {1};
\draw[fill=cube3,draw=black!87,line join=round,line width=0.674pt] (396,236) -- (396,256) -- (413.321,246) -- (413.321,226) -- cycle;
\node[anchor=center,inner sep=0pt,outer sep=0pt,text=black!93,font={\sffamily\fontsize{8}{9.6}\selectfont}] at (404.66,241) {3};
\draw[fill=cube2,draw=black!87,line join=round,line width=0.674pt] (396,236) -- (413.321,226) -- (396,216) -- (378.679,226) -- cycle;
\node[anchor=center,inner sep=0pt,outer sep=0pt,text=black!93,font={\sffamily\fontsize{8}{9.6}\selectfont}] at (396,226) {2};
\node[anchor=center,inner sep=0pt,outer sep=0pt,text=black!93,font={\sffamily\fontsize{9}{10.8}\selectfont}] at (396,271) {132};
\draw[fill=cube1,draw=black!87,line join=round,line width=0.674pt] (450,236) -- (432.679,226) -- (432.679,246) -- (450,256) -- cycle;
\node[anchor=center,inner sep=0pt,outer sep=0pt,text=black!93,font={\sffamily\fontsize{8}{9.6}\selectfont}] at (441.34,241) {1};
\draw[fill=cube4,draw=black!87,line join=round,line width=0.674pt] (450,236) -- (450,256) -- (467.321,246) -- (467.321,226) -- cycle;
\node[anchor=center,inner sep=0pt,outer sep=0pt,text=black!93,font={\sffamily\fontsize{8}{9.6}\selectfont}] at (458.66,241) {4};
\draw[fill=cube3,draw=black!87,line join=round,line width=0.674pt] (450,236) -- (467.321,226) -- (450,216) -- (432.679,226) -- cycle;
\node[anchor=center,inner sep=0pt,outer sep=0pt,text=black!93,font={\sffamily\fontsize{8}{9.6}\selectfont}] at (450,226) {3};
\node[anchor=center,inner sep=0pt,outer sep=0pt,text=black!93,font={\sffamily\fontsize{9}{10.8}\selectfont}] at (450,271) {143};
\draw[fill=cube0!30,draw=black!40,line join=round,line width=0.674pt] (504,236) -- (486.679,226) -- (486.679,246) -- (504,256) -- cycle;
\node[anchor=center,inner sep=0pt,outer sep=0pt,text=black!60,font={\sffamily\fontsize{8}{9.6}\selectfont}] at (495.34,241) {0};
\draw[fill=cube2!30,draw=black!40,line join=round,line width=0.674pt] (504,236) -- (504,256) -- (521.321,246) -- (521.321,226) -- cycle;
\node[anchor=center,inner sep=0pt,outer sep=0pt,text=black!60,font={\sffamily\fontsize{8}{9.6}\selectfont}] at (512.66,241) {2};
\draw[fill=cube3!30,draw=black!40,line join=round,line width=0.674pt] (504,236) -- (521.321,226) -- (504,216) -- (486.679,226) -- cycle;
\node[anchor=center,inner sep=0pt,outer sep=0pt,text=black!60,font={\sffamily\fontsize{8}{9.6}\selectfont}] at (504,226) {3};
\node[anchor=center,inner sep=0pt,outer sep=0pt,text=black!60,font={\sffamily\fontsize{9}{10.8}\selectfont}] at (504,271) {023};
\draw[fill=cube0!30,draw=black!40,line join=round,line width=0.674pt] (558,236) -- (540.679,226) -- (540.679,246) -- (558,256) -- cycle;
\node[anchor=center,inner sep=0pt,outer sep=0pt,text=black!60,font={\sffamily\fontsize{8}{9.6}\selectfont}] at (549.34,241) {0};
\draw[fill=cube4!30,draw=black!40,line join=round,line width=0.674pt] (558,236) -- (558,256) -- (575.321,246) -- (575.321,226) -- cycle;
\node[anchor=center,inner sep=0pt,outer sep=0pt,text=black!60,font={\sffamily\fontsize{8}{9.6}\selectfont}] at (566.66,241) {4};
\draw[fill=cube2!30,draw=black!40,line join=round,line width=0.674pt] (558,236) -- (575.321,226) -- (558,216) -- (540.679,226) -- cycle;
\node[anchor=center,inner sep=0pt,outer sep=0pt,text=black!60,font={\sffamily\fontsize{8}{9.6}\selectfont}] at (558,226) {2};
\node[anchor=center,inner sep=0pt,outer sep=0pt,text=black!60,font={\sffamily\fontsize{9}{10.8}\selectfont}] at (558,271) {042};
\draw[draw=black!87,line width=0.787pt] (158,284) -- (158,290) -- (472,290) -- (472,284);
\node[anchor=center,inner sep=0pt,outer sep=0pt,text=black!93,font={\sffamily\fontsize{9}{10.8}\selectfont\bfseries}] at (315,308) {6 private};
\node[anchor=west,inner sep=0pt,outer sep=0pt,text=black!93,font={\sffamily\fontsize{10}{12}\selectfont}] at (604,226) {6 corners need B;};
\node[anchor=west,inner sep=0pt,outer sep=0pt,text=black!93,font={\sffamily\fontsize{10}{12}\selectfont\bfseries}] at (604,252) {only 5 copies.};
\draw[draw=black!20,line width=0.674pt] (16,320) -- (784,320);
\draw[fill=cube4,draw=black!87,line join=round,line width=0.674pt] (48,334) -- (72,334) -- (72,358) -- (48,358) -- cycle;
\node[anchor=center,inner sep=0pt,outer sep=0pt,text=black!93,font={\sffamily\fontsize{10}{12}\selectfont}] at (60,346) {4};
\draw[fill=cube2,draw=black!87,line join=round,line width=0.674pt] (24,358) -- (48,358) -- (48,382) -- (24,382) -- cycle;
\node[anchor=center,inner sep=0pt,outer sep=0pt,text=black!93,font={\sffamily\fontsize{10}{12}\selectfont}] at (36,370) {2};
\draw[fill=cube0,draw=black!87,line join=round,line width=0.674pt] (48,358) -- (72,358) -- (72,382) -- (48,382) -- cycle;
\node[anchor=center,inner sep=0pt,outer sep=0pt,text=black!93,font={\sffamily\fontsize{10}{12}\selectfont}] at (60,370) {0};
\draw[fill=cube1,draw=black!87,line join=round,line width=0.674pt] (72,358) -- (96,358) -- (96,382) -- (72,382) -- cycle;
\node[anchor=center,inner sep=0pt,outer sep=0pt,text=black!93,font={\sffamily\fontsize{10}{12}\selectfont}] at (84,370) {1};
\draw[fill=cube3,draw=black!87,line join=round,line width=0.674pt] (96,358) -- (120,358) -- (120,382) -- (96,382) -- cycle;
\node[anchor=center,inner sep=0pt,outer sep=0pt,text=black!93,font={\sffamily\fontsize{10}{12}\selectfont}] at (108,370) {3};
\draw[fill=cube1,draw=black!87,line join=round,line width=0.674pt] (48,382) -- (72,382) -- (72,406) -- (48,406) -- cycle;
\node[anchor=center,inner sep=0pt,outer sep=0pt,text=black!93,font={\sffamily\fontsize{10}{12}\selectfont}] at (60,394) {1};
\node[anchor=center,inner sep=0pt,outer sep=0pt,text=black!93,font={\sffamily\fontsize{10}{12}\selectfont\bfseries}] at (72,428) {5 \(\times\) C};
\node[anchor=center,inner sep=0pt,outer sep=0pt,text=black!93,font={\sffamily\fontsize{9}{10.8}\selectfont}] at (72,450) {031241};
\draw[fill=cube0,draw=black!87,line join=round,line width=0.674pt] (180,376) -- (162.679,366) -- (162.679,386) -- (180,396) -- cycle;
\node[anchor=center,inner sep=0pt,outer sep=0pt,text=black!93,font={\sffamily\fontsize{8}{9.6}\selectfont}] at (171.34,381) {0};
\draw[fill=cube1,draw=black!87,line join=round,line width=0.674pt] (180,376) -- (180,396) -- (197.321,386) -- (197.321,366) -- cycle;
\node[anchor=center,inner sep=0pt,outer sep=0pt,text=black!93,font={\sffamily\fontsize{8}{9.6}\selectfont}] at (188.66,381) {1};
\draw[fill=cube1,draw=black!87,line join=round,line width=0.674pt] (180,376) -- (197.321,366) -- (180,356) -- (162.679,366) -- cycle;
\node[anchor=center,inner sep=0pt,outer sep=0pt,text=black!93,font={\sffamily\fontsize{8}{9.6}\selectfont}] at (180,366) {1};
\node[anchor=center,inner sep=0pt,outer sep=0pt,text=black!93,font={\sffamily\fontsize{9}{10.8}\selectfont}] at (180,411) {011};
\draw[fill=cube0,draw=black!87,line join=round,line width=0.674pt] (234,376) -- (216.679,366) -- (216.679,386) -- (234,396) -- cycle;
\node[anchor=center,inner sep=0pt,outer sep=0pt,text=black!93,font={\sffamily\fontsize{8}{9.6}\selectfont}] at (225.34,381) {0};
\draw[fill=cube2,draw=black!87,line join=round,line width=0.674pt] (234,376) -- (234,396) -- (251.321,386) -- (251.321,366) -- cycle;
\node[anchor=center,inner sep=0pt,outer sep=0pt,text=black!93,font={\sffamily\fontsize{8}{9.6}\selectfont}] at (242.66,381) {2};
\draw[fill=cube1,draw=black!87,line join=round,line width=0.674pt] (234,376) -- (251.321,366) -- (234,356) -- (216.679,366) -- cycle;
\node[anchor=center,inner sep=0pt,outer sep=0pt,text=black!93,font={\sffamily\fontsize{8}{9.6}\selectfont}] at (234,366) {1};
\node[anchor=center,inner sep=0pt,outer sep=0pt,text=black!93,font={\sffamily\fontsize{9}{10.8}\selectfont}] at (234,411) {021};
\draw[fill=cube1,draw=black!87,line join=round,line width=0.674pt] (288,376) -- (270.679,366) -- (270.679,386) -- (288,396) -- cycle;
\node[anchor=center,inner sep=0pt,outer sep=0pt,text=black!93,font={\sffamily\fontsize{8}{9.6}\selectfont}] at (279.34,381) {1};
\draw[fill=cube1,draw=black!87,line join=round,line width=0.674pt] (288,376) -- (288,396) -- (305.321,386) -- (305.321,366) -- cycle;
\node[anchor=center,inner sep=0pt,outer sep=0pt,text=black!93,font={\sffamily\fontsize{8}{9.6}\selectfont}] at (296.66,381) {1};
\draw[fill=cube3,draw=black!87,line join=round,line width=0.674pt] (288,376) -- (305.321,366) -- (288,356) -- (270.679,366) -- cycle;
\node[anchor=center,inner sep=0pt,outer sep=0pt,text=black!93,font={\sffamily\fontsize{8}{9.6}\selectfont}] at (288,366) {3};
\node[anchor=center,inner sep=0pt,outer sep=0pt,text=black!93,font={\sffamily\fontsize{9}{10.8}\selectfont}] at (288,411) {113};
\draw[fill=cube1,draw=black!87,line join=round,line width=0.674pt] (342,376) -- (324.679,366) -- (324.679,386) -- (342,396) -- cycle;
\node[anchor=center,inner sep=0pt,outer sep=0pt,text=black!93,font={\sffamily\fontsize{8}{9.6}\selectfont}] at (333.34,381) {1};
\draw[fill=cube2,draw=black!87,line join=round,line width=0.674pt] (342,376) -- (342,396) -- (359.321,386) -- (359.321,366) -- cycle;
\node[anchor=center,inner sep=0pt,outer sep=0pt,text=black!93,font={\sffamily\fontsize{8}{9.6}\selectfont}] at (350.66,381) {2};
\draw[fill=cube3,draw=black!87,line join=round,line width=0.674pt] (342,376) -- (359.321,366) -- (342,356) -- (324.679,366) -- cycle;
\node[anchor=center,inner sep=0pt,outer sep=0pt,text=black!93,font={\sffamily\fontsize{8}{9.6}\selectfont}] at (342,366) {3};
\node[anchor=center,inner sep=0pt,outer sep=0pt,text=black!93,font={\sffamily\fontsize{9}{10.8}\selectfont}] at (342,411) {123};
\draw[fill=cube1,draw=black!87,line join=round,line width=0.674pt] (396,376) -- (378.679,366) -- (378.679,386) -- (396,396) -- cycle;
\node[anchor=center,inner sep=0pt,outer sep=0pt,text=black!93,font={\sffamily\fontsize{8}{9.6}\selectfont}] at (387.34,381) {1};
\draw[fill=cube3,draw=black!87,line join=round,line width=0.674pt] (396,376) -- (396,396) -- (413.321,386) -- (413.321,366) -- cycle;
\node[anchor=center,inner sep=0pt,outer sep=0pt,text=black!93,font={\sffamily\fontsize{8}{9.6}\selectfont}] at (404.66,381) {3};
\draw[fill=cube4,draw=black!87,line join=round,line width=0.674pt] (396,376) -- (413.321,366) -- (396,356) -- (378.679,366) -- cycle;
\node[anchor=center,inner sep=0pt,outer sep=0pt,text=black!93,font={\sffamily\fontsize{8}{9.6}\selectfont}] at (396,366) {4};
\node[anchor=center,inner sep=0pt,outer sep=0pt,text=black!93,font={\sffamily\fontsize{9}{10.8}\selectfont}] at (396,411) {134};
\draw[fill=cube2,draw=black!87,line join=round,line width=0.674pt] (450,376) -- (432.679,366) -- (432.679,386) -- (450,396) -- cycle;
\node[anchor=center,inner sep=0pt,outer sep=0pt,text=black!93,font={\sffamily\fontsize{8}{9.6}\selectfont}] at (441.34,381) {2};
\draw[fill=cube4,draw=black!87,line join=round,line width=0.674pt] (450,376) -- (450,396) -- (467.321,386) -- (467.321,366) -- cycle;
\node[anchor=center,inner sep=0pt,outer sep=0pt,text=black!93,font={\sffamily\fontsize{8}{9.6}\selectfont}] at (458.66,381) {4};
\draw[fill=cube3,draw=black!87,line join=round,line width=0.674pt] (450,376) -- (467.321,366) -- (450,356) -- (432.679,366) -- cycle;
\node[anchor=center,inner sep=0pt,outer sep=0pt,text=black!93,font={\sffamily\fontsize{8}{9.6}\selectfont}] at (450,366) {3};
\node[anchor=center,inner sep=0pt,outer sep=0pt,text=black!93,font={\sffamily\fontsize{9}{10.8}\selectfont}] at (450,411) {243};
\draw[fill=cube0!30,draw=black!40,line join=round,line width=0.674pt] (504,376) -- (486.679,366) -- (486.679,386) -- (504,396) -- cycle;
\node[anchor=center,inner sep=0pt,outer sep=0pt,text=black!60,font={\sffamily\fontsize{8}{9.6}\selectfont}] at (495.34,381) {0};
\draw[fill=cube1!30,draw=black!40,line join=round,line width=0.674pt] (504,376) -- (504,396) -- (521.321,386) -- (521.321,366) -- cycle;
\node[anchor=center,inner sep=0pt,outer sep=0pt,text=black!60,font={\sffamily\fontsize{8}{9.6}\selectfont}] at (512.66,381) {1};
\draw[fill=cube4!30,draw=black!40,line join=round,line width=0.674pt] (504,376) -- (521.321,366) -- (504,356) -- (486.679,366) -- cycle;
\node[anchor=center,inner sep=0pt,outer sep=0pt,text=black!60,font={\sffamily\fontsize{8}{9.6}\selectfont}] at (504,366) {4};
\node[anchor=center,inner sep=0pt,outer sep=0pt,text=black!60,font={\sffamily\fontsize{9}{10.8}\selectfont}] at (504,411) {014};
\draw[fill=cube0!30,draw=black!40,line join=round,line width=0.674pt] (558,376) -- (540.679,366) -- (540.679,386) -- (558,396) -- cycle;
\node[anchor=center,inner sep=0pt,outer sep=0pt,text=black!60,font={\sffamily\fontsize{8}{9.6}\selectfont}] at (549.34,381) {0};
\draw[fill=cube4!30,draw=black!40,line join=round,line width=0.674pt] (558,376) -- (558,396) -- (575.321,386) -- (575.321,366) -- cycle;
\node[anchor=center,inner sep=0pt,outer sep=0pt,text=black!60,font={\sffamily\fontsize{8}{9.6}\selectfont}] at (566.66,381) {4};
\draw[fill=cube2!30,draw=black!40,line join=round,line width=0.674pt] (558,376) -- (575.321,366) -- (558,356) -- (540.679,366) -- cycle;
\node[anchor=center,inner sep=0pt,outer sep=0pt,text=black!60,font={\sffamily\fontsize{8}{9.6}\selectfont}] at (558,366) {2};
\node[anchor=center,inner sep=0pt,outer sep=0pt,text=black!60,font={\sffamily\fontsize{9}{10.8}\selectfont}] at (558,411) {042};
\draw[draw=black!87,line width=0.787pt] (158,424) -- (158,430) -- (472,430) -- (472,424);
\node[anchor=center,inner sep=0pt,outer sep=0pt,text=black!93,font={\sffamily\fontsize{9}{10.8}\selectfont\bfseries}] at (315,448) {6 private};
\node[anchor=west,inner sep=0pt,outer sep=0pt,text=black!93,font={\sffamily\fontsize{10}{12}\selectfont}] at (604,366) {6 corners need C;};
\node[anchor=west,inner sep=0pt,outer sep=0pt,text=black!93,font={\sffamily\fontsize{10}{12}\selectfont\bfseries}] at (604,392) {only 5 copies.};
\draw[draw=black!20,line width=0.674pt] (16,466) -- (784,466);
\end{tikzpicture}

\caption{The fourteen-cube obstruction $4A+5B+5C$. Each net has $+x$ at its
centre. Corner types are read counterclockwise from outside; the
full-colour types occur on no other variety. A target of variety $X$ has
these private positions, and only copies of $X$ fit them.}
\label{fig:obstruction}
\end{figure}

\section{Reduction to finite problems}\label{sec:cases}
Eight copies of any variety solve that variety, so a solution-free collection
has every multiplicity at most seven. Suppose such a collection has size 15.

\begin{lemma}[Anchored Hall cover]\label{lem:anchor}
Suppose a collection contains $m$ copies of a variety $A$ and does not solve
$A$. There is a set $K$ of its target positions such that
\begin{equation}\label{eq:anchor}
 |K|\le 7-m,\qquad
 \#\{\text{other copies fitting some position outside }K\}
       \le 7-m-|K|.
\end{equation}
Conversely, such a set certifies failure of target $A$.
\end{lemma}
\begin{proof}
Apply~\eqref{eq:hall} to $A$ and put $K=D^c$. Every copy of $A$ is a
neighbour of the nonempty set $D$. Subtract these $m$ copies from the bound
$|D|-1=7-|K|$. The converse is the same Hall inequality read backwards.
\end{proof}

First suppose the maximum multiplicity $m$ is at least two. Choose an
adjacent-repeat anchor attaining $m$ if possible; otherwise choose an
opposite-repeat anchor attaining $m$. In the latter case all adjacent-repeat
multiplicities are at most $m-1$. Relabel the colours to send the anchor to
a fixed representative of its orbit.

For each cover $K$ in Lemma~\ref{lem:anchor}, put
\[
 S(K)=N_A(K^c)\setminus\{A\},\qquad b(K)=7-m-|K|.
\]
Write $S=S(K)$ and $b=b(K)$. The case has
$x_A=m$, $x(S)\le b$, and the multiplicity bounds just described.
Additional upper bounds use two-variety solutions. If $a$ copies of $A$
and $8-a$ copies of $V$ build a target, with $a\le m$, then
$x_V\le 7-a$ in a solution-free collection. These tests are finite:
the two varieties must cover all eight target positions, $A$ must fit at
least $a$, and $V$ at least $8-a$. They are also sufficient: at most $a$
positions are unavailable to $V$,
and all of them are available to $A$. Assign these to $A$, extend the
assignment to $a$ positions that $A$ fits, and assign the remaining
$8-a$ positions to $V$.

Colour permutations and simultaneous reflection of the entire collection
preserve solvability. First identify equal pairs $(S(K),b(K))$, even when
they arise from inequivalent position subsets $K$; then retain one
representative per orbit of these pairs under the anchor stabiliser.
The counts below refer to these constraint pairs. Transporting a cover
also transports its constrained varieties, budget and multiplicities, so
no counterexample is lost. Global reflection transports the whole
collection; individual cubes still move only by proper rotations.
The repeated-variety cases have the following counts,
totalling 244.
\begin{center}
\begin{tabular}{lrrrrrr}
\toprule
Anchor / maximum multiplicity & 2 & 3 & 4 & 5 & 6 & 7\\
\midrule
Adjacent-repeat & 72 & 55 & 31 & 14 & 4 & 1\\
Opposite-repeat & 26 & 21 & 11 & 6 & 2 & 1\\
Total & 98 & 76 & 42 & 20 & 6 & 2\\
\bottomrule
\end{tabular}
\end{center}

For distinct collections a different choice of anchor reduces the count.
\begin{lemma}[Shared two-colour corner]\label{lem:shared}
Among eleven adjacent-repeat cubes, two share a two-colour corner type.
If a distinct collection contains such a pair and one member is chosen as
anchor, the anchor's two two-colour positions cannot be Hall-deficient.
\end{lemma}
\begin{proof}
Each adjacent-repeat cube has exactly two such types, $aab$ and $aac$, with
$b\ne c$. There are only $5\cdot4=20$ types $aab$. Eleven cubes contribute
22 incidences, so two cubes share a type. For the two characteristic positions
of the chosen anchor, the anchor itself and the other cube are two distinct
neighbours. Their neighbourhood therefore has size at least two.
\end{proof}
If an opposite-repeat variety is present, choose it as anchor; this gives
31 constraint orbits with every multiplicity at most one. Otherwise use
Lemma~\ref{lem:shared} to choose an adjacent-repeat anchor. Omit the cover
complementary to its two characteristic positions and set all opposite-repeat
capacities to zero. Its stabiliser preserves these positions, so the omission
is compatible with the orbit reduction. It removes one of the 83
adjacent-repeat constraint orbits, leaving 82 further cases.

\begin{proposition}\label{prop:cover}
Every solution-free fifteen-cube collection is carried by colour relabelling
and, if needed, global reflection into one of the 357 exported root cases.
\end{proposition}
\begin{proof}
The preceding arguments cover repeated and distinct collections. The finite
catalogue, two-variety bounds and coverage of all anchor-cover orbits are
checked from physical rotations by the accompanying audit. It regenerates
rotations from quarter turns, separately from the signed-axis permutations
used by the case generator, and checks every exported bound and every cover
orbit. The resulting counts are
\[
  244+31+82=357.
\]
The cases need not be disjoint; completeness of the cover is what is required.
\end{proof}

Each root specifies an anchor $A$ of multiplicity $m$, coordinate bounds
$\ell,u$, a set $S$ of non-anchor varieties and a budget $b$.
The bounds fix $\ell_A=u_A=m$, with $\ell_i=0$ for $i\ne A$.
A vector $x\in\mathbb Z_{\ge0}^{75}$ is \emph{admissible} for this root if
\begin{equation}\label{eq:admissible}
\begin{gathered}
 \ell\le x\le u,\qquad \sum_{i=1}^{75}x_i=15,\qquad x(S)\le b,\\
 \text{for every target }T\text{, there is }\varnothing\ne D\subseteq\{1,\ldots,8\}
 \text{ with }x(N_T(D))\le |D|-1.
\end{gathered}
\end{equation}
Thus admissibility includes failure of every target, in addition to the
anchor constraints.

The cover can be reduced without further search. If $S'\subseteq S$ and
$b'\ge b$, then $x(S)\le b$ implies $x(S')\le b'$ for nonnegative $x$.
Together with coordinatewise inclusion of the root boxes, this identifies
49 cases contained in other cases, leaving 308. Write $u$ for the root upper
bounds. In 29 retained cases, $b+u(S^c)<15$ already contradicts the prescribed
total, so only 279 cases
need further exclusion proofs. The exact map and retained case numbers are
in \path{data/q5-root-reduction.json}; Lean checks all containments and
arithmetic exclusions. The archive retains all 357 original cases.

\begin{lemma}[Hall-budget projection]\label{lem:projection}
Suppose a root has anchor multiplicity $m$, budget $x(S)\le b$, and total
$15$, with the anchor outside $S$. For $0\le s\le b$, put $G=15-b+s$
and suppose $G\ge m$. Every admissible vector $x$ has a vector $y\le x$
with the same anchor multiplicity, total $G$, budget $y(S)\le s$, and a
Hall deficiency for every target. It obeys the original upper bounds,
tightened to $\min(u_i,s)$ on $S$.
\end{lemma}
\begin{proof}
First retain only $\min(x(S),s)$ copies in $S$. This removes at most $b-s$
copies and leaves the anchor untouched, so the remaining total is at least
$G$. Delete further non-anchor copies until the total is $G$; this is possible
because $G\ge m$. Every coordinate decreases, so all upper bounds and Hall
deficiency inequalities are preserved. Nonnegativity and $y(S)\le s$ give
$y_i\le s$ for $i\in S$.
\end{proof}

For our anchored covers, $b=7-m-|K|$, so
\[
 G=m+8+|K|+s.
\]
Thus the reduced collection consists of the $m$ fixed anchor copies and
$8+|K|+s$ other copies; in particular, $G\ge m$ is automatic. Apply the lemma
after selecting and normalising the original root, retaining its bounds.
No new anchor selection is required for the smaller collection.
Of the 279 remaining roots, 277 admit an exclusion certificate with $s=0$. The other two, original
cases 68 and 69, use $s=1$ and $G=12$. Their $s=0$, $G=11$ problems actually
have solution-free collections, supplied and checked by physical matching in
the projected manifest. The projected totals range from 9 to 15. All 101 roots
whose original budget was zero retain their original certificates.

\section{Interval certificates}\label{sec:certificates}
For a projected problem, admissibility is defined by~\eqref{eq:admissible}
with tightened upper bounds and total $G$. In this section denote its
projected budget by $b$. We certify that no admissible vector exists.

Write a box as $\ell\le x\le u$. For a cut $x(S)\le r$, two necessary
conditions for feasibility in the box are
\begin{equation}\label{eq:possible}
  \ell(S)\le r,\qquad G\le r+u(S^c).
\end{equation}
The first follows from the lower bounds; the second follows from
$G=x(S)+x(S^c)$. When they hold, put
$\sigma=r-\ell(S)$ and $\gamma=r+u(S^c)-G$. Every satisfying vector obeys
\begin{equation}\label{eq:refine}
 x_i\le\min(u_i,\ell_i+\sigma)\quad(i\in S),\qquad
 x_i\ge\max(\ell_i,u_i-\gamma)\quad(i\notin S).
\end{equation}
For example, the upper bound follows by giving all other coordinates of $S$
their lower bounds; the lower bound follows by giving all other coordinates
of $S^c$ their upper bounds.

For a target, retain its cuts satisfying~\eqref{eq:possible}. If none remain,
no vector in the box can fail that target. Otherwise let $I$ and $U$ be the
intersection and union of their neighbourhoods, and let $\sigma_*,\gamma_*$ be the
largest corresponding $\sigma,\gamma$. Then all vectors in the box that fail this
target remain in the refined box
\begin{equation}\label{eq:summary}
 u'_i=\begin{cases}\min(u_i,\ell_i+\sigma_*)&i\in I,\\u_i&i\notin I,\end{cases}
 \qquad
 \ell'_i=\begin{cases}\max(\ell_i,u_i-\gamma_*)&i\notin U,\\\ell_i&i\in U.\end{cases}
\end{equation}
Indeed, a failing vector satisfies one of the retained cuts, and its bounds
in~\eqref{eq:refine} imply~\eqref{eq:summary}. The root budget is treated as
a family containing a single cut. The total $\sum_i x_i=G$ gives
\[
 \ell'_i=\max\!\left(\ell_i,G-\sum_{j\ne i}u_j\right),\qquad
 u'_i=\min\!\left(u_i,G-\sum_{j\ne i}\ell_j\right).
\]
A total-bound leaf closes exactly when $\sum_i\ell_i>G$ or
$\sum_i u_i<G$. A budget leaf closes when
$\ell(S)>b$ or $b+u(S^c)<G$; a target leaf closes when none of its
255 cuts satisfies~\eqref{eq:possible}. The checker thus has 77 closing
reasons: one total test, one budget test and 75 target tests. Refinements
must tighten at least one bound and leave
nonempty coordinate intervals; a split must satisfy $\ell_i\le v<u_i$.

A certificate is a finite tree of refinements and integer splits
$x_i\le v$ or $x_i\ge v+1$. Each leaf is excluded by total bounds or by a
cut family with no surviving cut. Induction on this tree proves that its
root contains no admissible vector. The separate C++ checker reads these
instructions and verifies the deductions; it performs no search and does
not use the search program's matching implementation.

\begin{proposition}[Certificate replay]\label{prop:replay}
The separate checker accepts an exclusion certificate for each of the
279 projected problems.
\end{proposition}
The projected archive contains 499,829 instructions, 95.69\% fewer than the
original archive. Both archives and the search computations are documented
in the repository.
Every archived header is checked against its exported root case. Eighty
negative controls are rejected, including all 77 possible direct closing
claims against a separate $G=11$ fixture, not a $G=15$ root problem.
The witness satisfies its total, fixed anchor count, coordinate bounds and
zero root budget, and fails all 75 targets by physical matching. The fixture
is \path{data/q5-distinct-known-control.txt}; its regeneration and witness
checks are in \path{scripts/check_five_colour_distinct.py}.

\begin{proof}[Proof of Theorem~\ref{thm:main}]
Lemma~\ref{lem:lower} gives $N_5\ge15$. A solution-free collection of fifteen
cubes would satisfy an exported root case by Proposition~\ref{prop:cover},
and would supply a Hall deficiency for every target. Containment and the initial
budget tests leave one of 279 roots. Lemma~\ref{lem:projection} supplies an
admissible vector in its projected problem, which Proposition~\ref{prop:replay}
and the sound interval deductions exclude. Hence $N_5\le15$.
\end{proof}

\section{Frames and complete cubes}\label{sec:frames}
The larger-cube argument depends on how many positions a source can miss.
The following form of Hall's theorem isolates the matching step.

\begin{lemma}[Hall criterion with a matched subset]\label{lem:partial-matching}
Consider $L$ target positions and $N\ge L$ source copies. Suppose every
source misses at most $d$ positions. If a set $P_0$ of positions can be
matched to distinct copies, and every position outside $P_0$ has at least
$d$ neighbours, then all $L$ positions can be matched.
\end{lemma}
\begin{proof}
For a position set $D\subseteq P_0$, the given matching supplies $|D|$
neighbours. If $D\nsubseteq P_0$ and $|D|\le d$, a position in
$D\setminus P_0$ alone has at least $d$ neighbours. If $|D|>d$, every
source neighbours $D$, since none can miss all its positions. In each case
Hall's inequality holds.
\end{proof}

\begin{lemma}[Frame compatibility]\label{lem:frame-geometry}
Put $q=n-2\ge1$. Every source variety fits at least $8q+2$ positions of
every five-colour $n$-frame target, and at least $9q+2$ if the target is
opposite-repeat. An adjacent-repeat source lacks exactly one bichromatic
edge type. An opposite-repeat source lacks exactly two, which are disjoint;
these two missing pairs determine its variety uniquely.
\end{lemma}
\begin{proof}
A missing bichromatic pair must occupy opposite singleton faces: if a colour
is repeated, a face of the other colour cannot be opposite both copies.
An adjacent-repeat variety has one opposite pair of singleton faces, and an
opposite-repeat variety has two. Two disjoint missing pairs specify two
opposite face pairs, leaving the fifth colour on the remaining opposite
faces, which determines a unique variety.

Write $c,e$ for the numbers of target corners and geometric edges fitted by
the source. Whenever either variety is opposite-repeat, $c\ge2$. To see
this, let that variety repeat $r$, with singleton opposite pairs
$\{a,b\}$ and $\{u,v\}$. It contains both cyclic orientations of each
$rxy$ with $x\in\{a,b\}$ and $y\in\{u,v\}$. In the other variety, if an
$r$-face has no adjacent $r$-face, its four surrounding faces contain at
least three distinct colours. Their cyclic order crosses between the two
pairs at least twice, giving two shared corner types. If the $r$-faces are
adjacent, the other neighbours of each form a path of three distinct
colours. Each path crosses between the pairs and gives a shared corner;
these corners have different types.

A bichromatic target edge type occurs at most twice, and a doubled type
contains the target's repeated colour. Two disjoint types therefore account
for at most three target edges. There is at most one monochromatic target
edge. These observations give the following bounds; A and O denote
adjacent-repeat and opposite-repeat varieties.
\begin{center}
\begin{tabular}{ccccc}
\toprule
Target & Source & $c$ & $e$ & $c+e$\\
\midrule
O & O & $\ge2$ & $\ge9$ & $\ge11$\\
O & A & $\ge2$ & $\ge10$ & $\ge12$\\
A & O & $\ge2$ & $\ge8$ & $\ge10$\\
A & A & $\ge0$ & $\ge9$ & $\ge10$\\
\bottomrule
\end{tabular}
\end{center}
Only the final entry needs an orientation argument. If an adjacent-repeat
source fits exactly nine edges of an adjacent-repeat target, it misses the
monochromatic edge and a doubled bichromatic type. Relabel colours and
rotate to put the target at $(0,1,0,2,3,4)$ and the missing type at $03$.
The source's $0$- and $3$-faces are opposite; its surrounding colours are
$1,2,4$, one repeated on adjacent faces. Suppress this adjacent repetition.
One of the two cyclic orders gives the shared target corner $123$; the
other gives $024$ and $041$. Thus $c\ge1$ when $e=9$.

Each geometric edge contributes $q$ positions. Since $e\ge8$ and $c+e\ge10$,
\[
 c+qe=(c+e)+(q-1)e\ge8q+2.
\]
For an opposite-repeat target, $e\ge9$ and $c+e\ge11$ instead give $9q+2$.
\end{proof}

\begin{theorem}\label{thm:frames}
For every integer $n\ge3$,
\[
 12n-16\le F_5(n)\le\max\{12n-16,\;11n-6\}.
\]
Consequently $F_5(n)=12n-16$ for $n\ge10$.
\end{theorem}
\begin{proof}
The lower bound is the frame's piece count. Put
\[
 q=n-2,\quad L=12q+8,\quad d=4q+6,\quad h=3q+6,
 \quad N=\max\{L,11q+16\}.
\]
By Lemma~\ref{lem:frame-geometry}, a source misses at most $d$ target
positions, or at most $h$ for an opposite-repeat target. Consider $N$ source
copies. Call a bichromatic pair \emph{rare} if at most $d-1=4q+5$ copies
possess that edge type.

If two distinct pairs are rare, at least
\[
 N-2(4q+5)\ge3q+6=h
\]
copies lack both. The lemma implies that these pairs are disjoint and all
these copies belong to one opposite-repeat variety $V$. Each of these copies
fits every position of $V$, so every position has at least $h$ neighbours.
Every source misses at most $h$ positions of this target.
Lemma~\ref{lem:partial-matching}, with $d=h$ and $P_0=\varnothing$,
therefore supplies a frame.

It remains to treat at most one rare pair. If there is none, reserve up to
$q$ cubes with an $aa$ edge for each of the five colours $a$, taking all if
fewer exist. These reserves are disjoint: a five-colour cube has only one
repeated colour. At least $N-5q\ge6q+16\ge15$ cubes remain, so
Theorem~\ref{thm:main} supplies a corner solution from them.

If exactly one pair $ab$ is rare, first remove every source possessing it.
At most $4q+5$ copies are removed. Every remaining source has $a,b$ on
opposite singleton faces, so only the other three colours can be repeated.
Reserve up to $q$ monochromatic-edge cubes for each of those colours.
At least
\[
 N-(4q+5)-3q\ge4q+11\ge15
\]
cubes remain, and again choose a corner solution. Its target cannot have
an $ab$ edge, since either endpoint corner would require a source
possessing that pair.

In both cases the chosen target uses no rare bichromatic type. If it has
a monochromatic edge $rr$, its corner solution contains an $rr$-capable
source outside the reserves. Thus the $rr$ reserve reached size $q$;
otherwise all such sources would have been removed. This reserve fills
the monochromatic edge. Hence the corners and the possible monochromatic
edge have a matching by distinct copies.

Return to all $N$ sources. Apply Lemma~\ref{lem:partial-matching} with
$P_0$ consisting of the corners and any monochromatic edge positions.
Each remaining position has at least $d$ neighbours because its type is
nonrare, and every source misses at most $d$ positions. This gives the
full frame. The final matching may rearrange the preliminary placements.

Substituting $q=n-2$ gives the stated bound. For $n\ge10$, its maximum is
$12n-16$, proving equality.
\end{proof}

\begin{theorem}\label{thm:complete}
For every integer $n\ge2$,
\[
 G_5(n)=\max\{n^3,15\}.
\]
In particular, every collection of $n^3$ five-colour cubes can all be
assembled into a complete $n$-cube for every $n\ge3$.
\end{theorem}
\begin{proof}
For $n=2$, a complete cube is its corner frame, so
Theorem~\ref{thm:main} gives $G_5(2)=15$. For $n\ge3$, the piece count
requires $G_5(n)\ge n^3$. Conversely, $n^3\ge12n-16$, since the difference
counts face-interior and interior positions, and
\[
 n^3-(11n-6)=(n-3)(n^2+3n-2)\ge0.
\]
Theorem~\ref{thm:frames} therefore supplies a frame from any $n^3$ sources.
Each remaining face-interior position prescribes just one colour, which
every unused cube has, and interior positions prescribe none. All remaining
cubes can fill these positions.
\end{proof}
Thus any 27 cubes suffice for a $3\times3\times3$ assembly without spares.
This does not determine the sharp frame value: it gives
$20\le F_5(3)\le27$. The sharp frame values for $3\le n\le9$ remain
undetermined here.

\section{Collections with few varieties}\label{sec:support}
Let $M(k)$ be the least number of cubes guaranteeing a corner solution when
the collection contains at most $k$ varieties.

\begin{lemma}\label{lem:disjoint-pair}
Two five-colour varieties with disjoint corner sets can build an
opposite-repeat target from four copies of each.
\end{lemma}
\begin{proof}
This is a small finite geometric check. Of the $\binom{75}{2}$ unordered
pairs, 330 have disjoint corner sets. Up to colour permutation and exchange
of the sources, they have the following five representatives. The first
source is $010234$ throughout; tuples use the face order of
Section~\ref{sec:model}. Each listed target has four corners supplied by
each source.
\begin{center}
\begin{tabular}{cc}
\toprule
Second source & Target\\
\midrule
$011243$ & $001234$\\
$012324$ & $012234$\\
$012334$ & $012234$\\
$012434$ & $012234$\\
$030412$ & $001234$\\
\bottomrule
\end{tabular}
\end{center}
The accompanying program \path{src/five_colour_extensions.cpp} regenerates
the varieties and checks the orbit coverage and all five witnesses.
\end{proof}

\begin{theorem}\label{thm:support}
For every integer $k\ge1$,
\[
 M(k)=\begin{cases}8&k=1,\\13&k=2,\\15&k\ge3.\end{cases}
\]
Consequently, three is the minimum number of varieties in a fourteen-cube
collection without a solution.
\end{theorem}
\begin{proof}
Eight identical copies build their own variety, and fewer than eight pieces
cannot form a corner solution. For two varieties, consider thirteen copies.
If either variety occurs at least eight times, it builds itself. Otherwise
the multiplicities are seven and six. If the varieties share a corner type,
use a minority copy at that corner of the majority target and seven majority
copies elsewhere. If their corner sets are disjoint, apply
Lemma~\ref{lem:disjoint-pair}.

For the lower bound take six copies each of
\[
 P=(0,1,0,2,3,4),\qquad Q=(0,1,1,3,2,4).
\]
Their corner sets share just $041$. Inspection of the 75 varieties shows
that only $P,Q$ have all corner types in their union. Each has seven private
positions, so neither can be built with six copies. The same geometric
checker verifies these assertions. Thus $M(2)=13$. For $k\ge3$,
Theorem~\ref{thm:main} supplies the upper bound, and the collection in
Lemma~\ref{lem:lower} supplies the lower bound.
\end{proof}

\begin{remark}[A family of fourteen-cube obstructions]
Keep $A,C$ from Section~\ref{sec:lower}, and put
\[
 B_1=(0,1,2,3,2,4),\quad B_2=(0,1,2,3,3,4),\quad B_3=(0,1,2,4,3,4).
\]
All 21 collections
\[
 4A+5C+x_1B_1+x_2B_2+x_3B_3,\qquad
 x_1+x_2+x_3=5,\quad x_i\in\mathbb Z_{\ge0},
\]
are unsolvable. The $B_i$ share the six trichromatic corner types
$023,034,042,124,132,143$. The combined corner union allows only
$A,C,B_1,B_2,B_3$ as targets, as a check of the catalogue confirms.
Target $A$ needs five $A$-copies, $C$ needs six $C$-copies, and any $B_i$
needs six copies from the $B$-group. The respective supplies are four,
five and five, regardless of their distribution within the $B$-group.
The original example is $x_1=5,x_2=x_3=0$.
\end{remark}

\section{Formalisation and reproducibility}\label{sec:formal}
The development uses Lean~4 and mathlib~\cite{Lean4,mathlib}, both at
\texttt{v4.35.0}. It proves the physical theorems, including the connection
between cube arrangements and the finite certificate problems.

Lean proves completeness and uniqueness of the 75 varieties and checks all
45,000 corner-compatibility bits using the shared corner--rotation theorem.
The lower bound uses 75 deficient corner sets: 72 singletons and three
sets of sizes 5, 6 and 6. Its theorem is
\path{FiveColour.Physical.lower_bound} in \path{FiveColour/LowerBound.lean}. This replaces the elementary
target classification in Lemma~\ref{lem:lower}; catalogue completeness and
rotation transport extend the result to every physical target.

The upper-bound formalisation follows the Hall reduction, symmetry
transport, anchor selection and shared-corner argument. Finite checks
verify capacities, cover representatives and the reduction from 357 roots
to 279. The deletion step is \path{FiveColour.Physical.root_projection}
in \path{FiveColour/RootProjection.lean}. The bridge to the archive, in
\path{FiveColour/ProjectedBridge.lean}, identifies the full admissibility
predicate: coordinate bounds, total, budget and all Hall-cut families.

For efficient replay, Lean proves soundness of a packed implementation of
the interval checker. For multiplicities $0\le x_i\le7$, index the varieties
by $0,\ldots,74$ and put
\[
 P(x)=\sum_{i=0}^{74}x_i B^i,\qquad
 W(S)=\sum_{i\in S}B^{74-i},\qquad B=1024.
\]
Then
\[
 x(S)=\left\lfloor\frac{P(x)W(S)}{B^{74}}\right\rfloor\bmod B.
\]
Every coefficient in the product is at most $75\cdot7<1024$, so no carry
can affect this coefficient extraction. Lean proves the identity, decoding,
coordinate updates and equivalence with the integer inequalities, and
checks all 19,125 neighbourhood weights.

The 279 certificates are split into 4,786 chunks of at most 128
instructions. Composition requires exact boundary-vector agreement and a
proof for every referenced box; the exporter is untrusted. All chunks and
their final assembly have passed kernel checking. The case mapping and
certificate hashes are in \path{data/q5-projected-root-manifest.json}.
The generated module \path{FiveColour/ProjectedChunks/Threshold.lean}
proves the unconditional theorems \path{threshold15} and
\path{sharpThreshold15}. Both lie in namespace
\path{FiveColour.Packed.ProjectedChunks}.

The larger-cube argument uses the general Hall lemmas of
\path{FiveColour/FiniteHall.lean}, including
Lemma~\ref{lem:partial-matching}, and their application in
\path{FiveColour/RareFrames.lean}. Lean checks the geometric bounds for all
5,625 source--target pairs and the variety determined by two missing pairs.
Edge bits are checked against face adjacency; corner counts reuse the
verified incidence table. The remaining argument applies to every
$q=n-2\ge1$. Frame completion in \path{FiveColour/FullCube.lean} adapts the
general construction from~\cite{FourToThree} without using its threshold.
The small-support proofs check pair witnesses directly, without reducing
them to five colour orbits, and use deficient corner sets for the two
obstruction constructions. They are in \path{FiveColour/SmallSupport.lean}
and \path{FiveColour/MixedObstructions.lean}.

The final statements in \path{Solution.lean}, in namespace
\path{FiveColour.Puzzle}, include \path{frame_bound}, \path{frame_threshold},
\path{full_cube}, \path{complete_cube_threshold} and \path{support_threshold}.
The 100 declarations in \path{Audit.lean}, the final replay theorems and the
seven statements in \path{ExtensionsAudit.lean} have audited dependencies
confined to \texttt{propext}, \texttt{Classical.choice} and
\texttt{Quot.sound}. No \texttt{native\_decide} is used. Reproduce the proof
and audits with:
\begin{verbatim}
make check WORKERS=2
make projected-replay-data WORKERS=8
make projected-replay WORKERS=8
make check-extensions
\end{verbatim}
The module guide is \path{FiveColour/README.md}.

A Mathlib-only comparator challenge states the corner problem geometrically.
Its equivalence with the development's model is proved by
\path{FiveColour.Puzzle.hasSolution_iff} in
\path{FiveColour/ComparatorBridge.lean}. The solution module applies the
completed threshold theorem. A full external comparator pass remains
unconfirmed; it is separate from the completed Lean proof.

\section{Remaining questions}
The frame values for $3\le n\le9$ remain undetermined here. In particular,
does $F_5(3)=20$, so that a $3$-frame never needs spare cubes? A second
question is to classify all fourteen-cube corner obstructions. The support
bound and the family in Section~\ref{sec:support} give initial constraints
and examples.

\section*{AI disclosure and review status}
GPT-6-Astra (OpenAI), used through Codex, developed the arguments, search and
verification programs, Lean formalisation and manuscript under the human
author's direction. The author initiated and directed the research. Author
review and submission approval of this working manuscript are pending.
The inherited geometry has its own provenance in the pinned dependency.
Automated checks and AI review are not independent human peer review.

\begin{thebibliography}{9}
\raggedright
\bibitem{MacMahon} P.~A.~MacMahon,
\emph{New Mathematical Pastimes}, Cambridge University Press (1921).
\href{https://archive.org/details/newmathematicalp00macmuoft}{Digitised edition}.
\bibitem{BES} E.~Berkove, A.~Eickemeyer and M.~Schulman,
\emph{The 4-Color Cubes Puzzle}, Integers \textbf{18} (2018), A68.
\href{https://doi.org/10.5281/zenodo.10680216}{doi:\nolinkurl{10.5281/zenodo.10680216}}.
\href{https://math.colgate.edu/~integers/s68/s68.pdf}{Journal PDF}.
\bibitem{Hall} P.~Hall, \emph{On Representatives of Subsets},
Journal of the London Mathematical Society \textbf{s1-10}(1) (1935), 26--30.
\href{https://doi.org/10.1112/jlms/s1-10.37.26}{doi:\nolinkurl{10.1112/jlms/s1-10.37.26}}.
\bibitem{FourToThree} J.~Davies,
\emph{Three-colour solutions from four-colour cubes}, manuscript (2026).
\href{https://github.com/jasondavies/four-to-three-colour-cubes}{Accompanying repository}.
\bibitem{ColourCubes} J.~Davies,
\emph{Colour cubes core}, Lean library (2026).
\href{https://github.com/jasondavies/colour-cubes-core/tree/3164adbd2ee16d2721caadc1db67c1498e5d67e0}{Pinned repository revision}.
\bibitem{Rosenfeld} M.~Rosenfeld,
\emph{MacMahon Cubes and Moshe's Insanity, Puzzle 1},
Geombinatorics Quarterly \textbf{36}(2) (October 2026), starts at p.~78.
\href{https://geombina.uccs.edu/past-issues/volume-xxxvi}{Journal listing};
full text not examined.
\bibitem{Lean4} L.~de Moura and S.~Ullrich,
\emph{The Lean 4 Theorem Prover and Programming Language},
in Automated Deduction---CADE~28, LNCS \textbf{12699}, Springer (2021), 625--635.
\href{https://doi.org/10.1007/978-3-030-79876-5_37}{doi:\nolinkurl{10.1007/978-3-030-79876-5_37}}.
\bibitem{mathlib} The mathlib Community, \emph{The Lean Mathematical Library},
in Proceedings of CPP~2020, ACM (2020), 367--381.
\href{https://doi.org/10.1145/3372885.3373824}{doi:\nolinkurl{10.1145/3372885.3373824}}.
\href{https://arxiv.org/abs/1910.09336}{arXiv:1910.09336}.
\end{thebibliography}
\end{document}
